\documentclass{article}
\usepackage[T1]{fontenc}
\usepackage{lmodern}
\usepackage{graphicx}
\usepackage{amsmath,amssymb}
\usepackage[a4paper,margin=2cm]{geometry}
\usepackage[absolute,overlay]{textpos}
\setlength{\TPHorizModule}{1mm}
\setlength{\TPVertModule}{1mm}
\usepackage[indonesian]{babel}
\usepackage{hyperref}
\setlength{\emergencystretch}{2em}
% unit: Halaman 77

\begin{document}

% === Halaman 77 (llm) ===

\section*{Kelebihan Trivium}

\begin{enumerate}
    \item Trivium cocok untuk \textit{lightweight cryptography}, karena \textit{state} relatif kecil (hanya 288 bit)
    \item \textit{Hardware efficiency}. Sederhana jika diimplementasikan secara \textit{hardware}, karena struktunta sederhana (terutama terdiri dari \textit{flip-flop/register}, gerbang AND dan OR)
\end{enumerate}

\end{document}
